% ======================================================================
%  Preambule commun - Sujets Bac Serie A (Mathematiques)
%  Style sobre : noir et blanc, sans couleurs, sans filigrane,
%  sans en-tetes ni pieds de page decoratifs.
% ======================================================================
\usepackage[utf8]{inputenc}
\usepackage[T1]{fontenc}
\usepackage{lmodern}
\usepackage{textcomp}
\usepackage{amsmath,amssymb}
\usepackage{enumitem}
\usepackage{array}
\usepackage[a4paper,margin=2.3cm]{geometry}
\usepackage{fancyhdr}
\usepackage{tikz}

% Symbole degre robuste + justification souple
\DeclareUnicodeCharacter{00B0}{\ensuremath{^\circ}}
\tolerance=2000
\emergencystretch=2em
\frenchspacing

% En-tete / pied de page minimalistes (numero de page seul)
\pagestyle{fancy}
\fancyhf{}
\renewcommand{\headrulewidth}{0pt}
\fancyfoot[C]{\thepage}

% Macros mathematiques
\newcommand{\R}{\mathbb{R}}
\newcommand{\N}{\mathbb{N}}
\newcommand{\Z}{\mathbb{Z}}
\newcommand{\vi}{\vec{\imath}}
\newcommand{\vj}{\vec{\jmath}}

% Titres de blocs
\newcommand{\exo}[2]{\bigskip\noindent{\large\bfseries Exercice #1}\hfill\textit{#2}\par\nopagebreak\smallskip}
\newcommand{\partie}[1]{\medskip\noindent{\bfseries Partie #1}\par\nopagebreak\smallskip}
\newcommand{\entetesujet}[1]{\begin{center}{\Large\bfseries #1}\end{center}\vspace{0.3em}\hrule\vspace{1.1em}}

% Questions numerotees : niveau 1 = chiffres gras, niveau 2 = a. b. c.
\newlist{questions}{enumerate}{1}
\setlist[questions]{label=\textbf{\arabic*.},leftmargin=1.8em,itemsep=3pt,topsep=3pt}
\newlist{sousq}{enumerate}{1}
\setlist[sousq]{label=\textbf{\alph*.},leftmargin=1.8em,itemsep=2pt,topsep=2pt}
